% Ad Atticum 13.28 --- headnote
% ad-atticum-13-28, 1 June 45 BC, Tusculanum

\textit{Cicero to Atticus, written at the Tusculan villa on 1 June
45 BC --- Perseus dateline \textit{Scr. in Tusculano Vii. K. Iun.
a. 709 (45)}. The sequel to 13.27 and the closing letter of the
late-May Tusculan run before the early-June continuation: the
dedication-shuffle around Varro and the \textit{Academica}
revision proceed in the background, but the foreground is again
the abandoned letter of advice to Caesar, on which Cicero now
declares himself defeated --- not by the indecency of flattering
the dictator (which, he admits in a sharp aside, ought to deter
him) but by the simple fact that nothing useful occurs to him.}

\textit{Section 2 is the heart: Cicero contrasts his blank
draftsman's slate with the schoolroom panegyrists who could
exhort Alexander to glory because their young king was ablaze
for it, and confesses that what he ``carved out of an oak'' was
criticized precisely for the few decent strokes in it. Section
3 sharpens the historical analogy: even Aristotle's pupil
turned tyrant once the kingship took hold, so what hope of
restraint in this \textit{contubernalis} of Quirinus's
procession (a glance at the statue of Caesar carried in the
parade of the gods)? One Greek phrase: \textit{probl\=ema
Archid\=emou}, ``the puzzle of Archidemus,'' an unidentified
Stoic-dialectical reference that does duty here as the old
spur Cicero had felt to write the letter, now extinguished.
The autograph postscript at section 4 is a piece of marriage-
market gossip: Thalna's suit for Cornificia, blocked by the
women of the house on financial grounds, with the dowry
spelled out in figures (the Latin abbreviates the sum: 800{,}000
sesterces).}
